\documentclass[11pt]{article}
\usepackage[margin=1in]{geometry}
\usepackage[T1]{fontenc}
\usepackage[utf8]{inputenc}
\usepackage{setspace}
\usepackage{enumitem}
\usepackage{hyperref}
\hypersetup{colorlinks=true, allcolors=black}
\setlength{\parindent}{0pt}
\setlength{\parskip}{0.3em}
\setstretch{1.0}
\pagestyle{empty}

\begin{document}

\begin{center}
{\Large Statement on the Integration of Teaching and Research}\\[0.25em]
Ke Cheng Hsu \\[0.15em]
{\small Duke Kunshan University faculty application (AJO \#32544)}
\end{center}

\noindent\textbf{Research and research agenda.} My research lies in industrial organization and digital economics. I have two working papers and two works in progress. My work focuses mainly on product return policies, with one project on consumer search and another on dynamic oligopoly under demand uncertainty.

\begin{itemize}[leftmargin=1.6em,itemsep=0em,topsep=0.05em,parsep=0em]
\item \textbf{Designing Return Policy in E-Commerce} (job market paper). A return policy affects two margins, whether shoppers buy and whether they return the item after delivery. A stricter policy cuts wasteful returns, but lower return rates make a higher price more profitable for the seller and can price out valuable purchases. I estimate a structural model of purchase, insurance, and return decisions across six apparel categories. Pairing a discounted price with a moderate restocking fee balances both margins better than free returns or optional insurance.
\item \textbf{Product Design with Heterogeneous Search Costs} (with Pak Hung Au). We introduce search-cost heterogeneity into a sequential search model. Low-search-cost consumers seek niche products, while high-search-cost consumers buy standardized products. When search costs fall, more consumers switch to searching niche products, which can raise prices for both product types.
\item \textbf{Pass-Through and Incidence with Product Returns} (with Ce Bian and Liang Zhao, work in progress). When a higher price raises return probability, the seller earns a smaller effective margin and cuts supply. A tax can therefore lower the equilibrium price and raise social welfare by reducing returns.
\item \textbf{Too Much Unsold Stock?} (with Kohei Kawaguchi and Kei Kawai, work in progress). In the Japanese book market, publishers fix retail prices and bookstores return unsold copies to be pulped. Fixed prices encourage bookstores to stock more copies under uncertain demand, but prevent discounting when demand is low. We estimate a dynamic oligopoly model on monthly bookstore data to measure the waste and evaluate a deposit-refund tax.
\end{itemize}

My five-year agenda extends the study of product returns in three steps. First, I will study search, purchase, and return decisions jointly because search resolves match uncertainty only partially. Second, I will study how substitutes and complements affect returns. In bracketing, a shopper orders several sizes or colors of the same item, keeps one, and returns the rest because those variants are perfect substitutes. I will test whether encouraging complementary purchases can reduce these returns. Third, I will study how return policies on two-sided platforms shape buyer and seller participation and pricing.

\noindent\textbf{How research improves my teaching.} Doing research on everyday markets helps me explain economic mechanisms more clearly to students. In class, I want students to look past formulas and understand why an outcome occurs. Because my research studies familiar settings like online shopping and bookstores, I can turn abstract models into concrete examples. For example, in a microeconomics class, I use online purchases and returns to show how uncertainty shapes a buyer's decision and therefore the seller's incentive to set the price. A stricter return rule cuts wasteful returns, yet fewer returns make a higher price more profitable for the seller and shut out buyers who would have kept the item. Working through these mechanisms in research helps me break arguments into simple steps and answer student questions in plain language. On the Professor of the Practice track at Duke Kunshan University, I will focus on teaching and mentoring while keeping an active research agenda that refreshes my courses and supports Signature Work.

\end{document}
